% GENERATED FILE. Edit canonical lessons, then run npm run generate:books.
% canonical-source-hash: fnv1a64:d04f604ddd27202f
% canonical-lessons: FR-C380-norme, FR-C380-reglement, FR-C380-obligatoire, FR-C380-facultatif, FR-C380-obligation

\chapter{Les règles}
\label{ch:fr-les-regles}
\hlchaptermodality{380}

\begin{chapteropening}
\textbf{By the end of this chapter:} \emph{I can talk about rules: a standard, the regulations, what is compulsory or optional, and an obligation.}

Complete the last lesson of chapter 380: I can talk about rules: a standard, the regulations, what is compulsory or optional, and an obligation.
\end{chapteropening}

\section[la norme]{la norme — the norm, the standard}

\label{lesson:FR-C380-norme}



\begin{quote}
Before the new one: say the French for the doctorate, then the French for the grant.
\end{quote}

\subsection*{You'll want to know: la norme}
\textbf{la norme} — \textquotedblleft{}the norm, the standard\textquotedblright{}.

\begin{quote}
Ce casque est aux normes.
\end{quote}

\begin{quote}
\emph{This helmet meets the standards.}
\end{quote}

\emph{Être aux normes} — to meet the standards.

\begin{cousinweb}
\emph{norme} is a learned borrowing from Latin \emph{nōrma} \textquotedblleft{}a carpenter's square, a rule, a pattern\textquotedblright{} (EtymDB; Lewis \& Short).
\end{cousinweb}

\begin{grammarlens}[title={What you've built}]
One more word for telling, explaining and describing at B1.
\end{grammarlens}

\subsection*{Your turn}
\begin{itemize}
\raggedright
  \item \emph{Say it:} \emph{la norme}
  \item \emph{Say it:} \emph{la norme}, once more
  \item \emph{Say it:} a sentence of your own with \emph{la norme}
  \item \emph{Recall:} say the French for the doctorate, then the French for the grant, then say \emph{la norme} again
\end{itemize}

\begin{tcolorbox}[breakable,colback=teal!4,colframe=teal!35!black,title={Before you move on}]
What does la norme mean? (\textquotedblleft{}The norm, the standard\textquotedblright{}.) Say it once more.
\end{tcolorbox}
\section[le règlement]{le règlement — the regulations, the rules; the payment}

\label{lesson:FR-C380-reglement}



\begin{quote}
Before the new one: say the French for the grant, then the French for the norm.
\end{quote}

\subsection*{You'll want to know: le règlement}
\textbf{le règlement} — \textquotedblleft{}the regulations, the rules; the payment\textquotedblright{}.

\begin{quote}
Le règlement de l'école interdit les téléphones en classe.
\end{quote}

\begin{quote}
\emph{The school rules forbid phones in class.}
\end{quote}

It also means payment: \emph{règlement par carte}.

\begin{cousinweb}
\emph{règlement} is built from \emph{régler} \textquotedblleft{}to regulate, to settle\textquotedblright{} and \emph{-ment}. The verb rests on \emph{règle} \textquotedblleft{}rule\textquotedblright{}, from Latin \emph{rēgula} \textquotedblleft{}a ruler, a rule\textquotedblright{}, built on \emph{regere} \textquotedblleft{}to guide\textquotedblright{} (Lewis \& Short).
\end{cousinweb}

\begin{grammarlens}[title={What you've built}]
One more word for telling, explaining and describing at B1.
\end{grammarlens}

\subsection*{Your turn}
\begin{itemize}
\raggedright
  \item \emph{Say it:} \emph{le règlement}
  \item \emph{Say it:} \emph{le règlement}, once more
  \item \emph{Say it:} a sentence of your own with \emph{le règlement}
  \item \emph{Recall:} say the French for the grant, then the French for the norm, then say \emph{le règlement} again
\end{itemize}

\begin{tcolorbox}[breakable,colback=teal!4,colframe=teal!35!black,title={Before you move on}]
What does le règlement mean? (\textquotedblleft{}The regulations, the rules; the payment\textquotedblright{}.) Say it once more.
\end{tcolorbox}
\section[obligatoire]{obligatoire — compulsory, mandatory}

\label{lesson:FR-C380-obligatoire}



\begin{quote}
Before the new one: say the French for the norm, then the French for the regulations.
\end{quote}

\subsection*{You'll want to know: obligatoire}
\textbf{obligatoire} — \textquotedblleft{}compulsory, mandatory\textquotedblright{}.

\begin{quote}
Le port du casque est obligatoire sur le chantier.
\end{quote}

\begin{quote}
\emph{Wearing a helmet is compulsory on the building site.}
\end{quote}

One form for both genders.

\begin{cousinweb}
\emph{obligatoire} is a learned borrowing from Late Latin \emph{obligātōrius} \textquotedblleft{}binding\textquotedblright{} (EtymDB), from \emph{obligāre} \textquotedblleft{}to bind\textquotedblright{}, built on \emph{ligāre} \textquotedblleft{}to tie\textquotedblright{} (Lewis \& Short).
\end{cousinweb}

\begin{grammarlens}[title={What you've built}]
One more word for telling, explaining and describing at B1.
\end{grammarlens}

\subsection*{Your turn}
\begin{itemize}
\raggedright
  \item \emph{Say it:} \emph{obligatoire}
  \item \emph{Say it:} \emph{obligatoire}, once more
  \item \emph{Say it:} a sentence of your own with \emph{obligatoire}
  \item \emph{Recall:} say the French for the norm, then the French for the regulations, then say \emph{obligatoire} again
\end{itemize}

\begin{tcolorbox}[breakable,colback=teal!4,colframe=teal!35!black,title={Before you move on}]
What does obligatoire mean? (\textquotedblleft{}Compulsory, mandatory\textquotedblright{}.) Say it once more.
\end{tcolorbox}
\section[facultatif]{facultatif — optional}

\label{lesson:FR-C380-facultatif}



\begin{quote}
Before the new one: say the French for the regulations, then the French for compulsory.
\end{quote}

\subsection*{You'll want to know: facultatif}
\textbf{facultatif} — \textquotedblleft{}optional\textquotedblright{}.

\begin{quote}
Le cours de musique est facultatif.
\end{quote}

\begin{quote}
\emph{The music class is optional.}
\end{quote}

Feminine \emph{facultative}.

\begin{cousinweb}
\emph{facultatif} is built on \emph{faculté} \textquotedblleft{}power, option\textquotedblright{}, from Latin \emph{facultās} \textquotedblleft{}ability, power\textquotedblright{} (Lewis \& Short), from \emph{facilis} \textquotedblleft{}easy, doable\textquotedblright{}: you may do it if you choose.
\end{cousinweb}

\begin{grammarlens}[title={What you've built}]
One more word for telling, explaining and describing at B1.
\end{grammarlens}

\subsection*{Your turn}
\begin{itemize}
\raggedright
  \item \emph{Say it:} \emph{facultatif}
  \item \emph{Say it:} \emph{facultatif}, once more
  \item \emph{Say it:} a sentence of your own with \emph{facultatif}
  \item \emph{Recall:} say the French for the regulations, then the French for compulsory, then say \emph{facultatif} again
\end{itemize}

\begin{tcolorbox}[breakable,colback=teal!4,colframe=teal!35!black,title={Before you move on}]
What does facultatif mean? (\textquotedblleft{}Optional\textquotedblright{}.) Say it once more.
\end{tcolorbox}
\section[l'obligation]{l'obligation — the obligation, the duty}

\label{lesson:FR-C380-obligation}



\begin{quote}
Before the new one: say the French for compulsory, then the French for optional.
\end{quote}

\subsection*{You'll want to know: l'obligation}
\textbf{l'obligation} — \textquotedblleft{}the obligation, the duty\textquotedblright{}.

\begin{quote}
Tu n'as aucune obligation de venir.
\end{quote}

\begin{quote}
\emph{You are under no obligation to come.}
\end{quote}

Feminine: \emph{une obligation}.

\begin{cousinweb}
\emph{obligation} is a learned borrowing from Latin \emph{obligātiō} \textquotedblleft{}a binding, an engagement\textquotedblright{} (EtymDB), from \emph{obligāre} \textquotedblleft{}to bind\textquotedblright{}, built on \emph{ligāre} \textquotedblleft{}to tie\textquotedblright{} (Lewis \& Short).
\end{cousinweb}

\begin{grammarlens}[title={What you've built}]
One more word for telling, explaining and describing at B1.
\end{grammarlens}

\subsection*{Your turn}
\begin{itemize}
\raggedright
  \item \emph{Say it:} \emph{l'obligation}
  \item \emph{Say it:} \emph{l'obligation}, once more
  \item \emph{Say it:} a sentence of your own with \emph{l'obligation}
  \item \emph{Recall:} say the French for compulsory, then the French for optional, then say \emph{l'obligation} again
\end{itemize}

\begin{tcolorbox}[breakable,colback=teal!4,colframe=teal!35!black,title={Before you move on}]
What does l'obligation mean? (\textquotedblleft{}The obligation, the duty\textquotedblright{}.) Say it once more.
\end{tcolorbox}
